%% ============================================================
%% CHAPTER 9 — CONCLUSIONS
%% ============================================================
\chapter{Conclusions and Recommendations}
\label{ch:conclusions}

\section{Answers to the Research Questions}
\label{sec:answers}

\subsection{RQ1: Thermodynamic Performance}

Under consistent assumptions --- equal turbomachinery efficiencies, a
recuperator limited by both streams, and a 36\,\textdegree C sCO$_2$
compressor inlet consistent with Australian cooling water --- the steam
Rankine cycle has the highest thermal efficiency at every turbine inlet
temperature from 150 to 600\,\textdegree C (19.1--35.4\,\%), followed by
tCO$_2$ (13.3--30.1\,\%) and sCO$_2$ (5.3--30.1\,\%); the sCO$_2$ cycle is
unviable below 300\,\textdegree C because its compressor absorbs 45--80\,\%
of turbine work. Per unit of net work steam is also the most
exergy-efficient cycle at a common turbine inlet temperature; at realistic
pinch-limited boiler pressures tCO$_2$ leads at 200 and 350\,\textdegree C.
The tCO$_2$ heater is pinched at both terminals with the driving force
widest in the middle, so the pseudo-critical $c_p$ peak enlarges the mean
temperature difference rather than restricting it. Recuperation raises
efficiency by 3--10 points at $\varepsilon = 0.70$ but raises the gas outlet
temperature with it, so it adds only 8--18\,\% to the work recovered per
kilogram of gas. The tCO$_2$ baseline needs a 14\,\textdegree C condenser
that a chiller could supply only by consuming 63--134\,\% of the net
output, and the sCO$_2$ cycle loses two-fifths of its 200\,\textdegree C
efficiency if its compressor inlet warms from 32 to 40\,\textdegree C.

\subsection{RQ2: Equipment Sizing and Capital Cost}

CO$_2$ turbines are 15--26 times more compact than the steam turbine at the
same site and remain single-stage (VFR 2.0--3.1) where steam needs
multi-stage machines above 350\,\textdegree C; this advantage is robust.
Integrated with the site stream, the steam boiler must run at 0.16, 0.91 and
3.69\,MPa, and the tCO$_2$ cycle recovers the most power at 200 and
350\,\textdegree C (723 and 3418\,kW against 412 and 2639\,kW for steam)
because its cold heater inlet cools the gas furthest; sCO$_2$ recovers the
least at every site. At equal output steam is the cheapest plant at every
site (AUD\,4{,}444, 3{,}071 and 2{,}598/kW at 1\,MWe), tCO$_2$ costs 30--65\,\%
more and sCO$_2$ two to three times more, the latter because its back-work
ratio requires 2.8\,MW of turbine and 1.8\,MW of compressor per megawatt
delivered at 200\,\textdegree C. Two-point LMTD sizing overestimates the
tCO$_2$ heater area by 129--223\,\% when the $c_p$ peak lies inside it.

\subsection{RQ3: Australian Techno-Economic Analysis}

On a fixed-source basis the steam cycle gives the lowest LCOE at all three
sites (AUD\,96, 35 and 31/MWh) and the highest NPV at 350 and
500\,\textdegree C (AUD\,11.4\,M and 12.4\,M, discounted payback 4.2 and
3.2 years, IRR 28 and 36\,\%, 99.8--100\,\% Monte Carlo viability). The
tCO$_2$ cycle is the better investment at 200\,\textdegree C (NPV 1.71\,M
against 1.12\,M) because it recovers 75\,\% more power, but only where a
natural cold sink exists; both options there are marginal (payback 10--11
years, 68--72\,\% viability). The sCO$_2$ cycle is not viable at any site
(NPV negative at CS2 and CS1; 0.5--49\,\% viability). Cheaper heat exchangers
would bring tCO$_2$ level with steam at 200\,\textdegree C with an 11\,\%
reduction, but not at 350 or 500\,\textdegree C, and would not rescue sCO$_2$
anywhere.

\subsection{RQ4: Off-Design Performance}

The tCO$_2$ cycle is the most flexible (Marchionni score 0.78; 56--67\,\%
efficiency retention at 20\,\% load), steam second (0.73; 38--66\,\%, worst
at the low-pressure 200\,\textdegree C design) and sCO$_2$ third (0.71;
41--47\,\%). The sCO$_2$ efficiency rise at 80\,\% load at 200\,\textdegree C
is a pressure-ratio effect of sliding-pressure control, not a near-critical
effect.

\section{Original Contributions}
\label{sec:contributions}

\begin{enumerate}[leftmargin=1.8em,itemsep=5pt]
\item \textbf{A consistent three-cycle Australian benchmark.} All three
cycles modelled with the same property library, turbomachinery efficiencies,
recuperator definition, heater integration, cost correlations and economic
basis, at three real Australian sites, with every number regenerated from one
script. No published study does this for SRC, sCO$_2$ and tCO$_2$ together.
\item \textbf{Quantification of why the CO$_2$ advantage is
conditional.} The thesis isolates four boundary conditions --- recuperator
energy balance, heat-sink temperature, steam baseline and cycle layout ---
and quantifies each: a cold-side-only recuperator overstates tCO$_2$
efficiency by 9--10 points; a 32$\to$40\,\textdegree C compressor inlet
removes 42\,\% of sCO$_2$ efficiency at 200\,\textdegree C; a 14\,\textdegree C
condenser costs more to provide than the tCO$_2$ cycle produces at
200\,\textdegree C.
\item \textbf{Corrected heater integration for near-critical CO$_2$.} With
a duty-linear gas profile the tCO$_2$ heater exhibits a double pinch with a
mid-exchanger bulge of 42--65\,K, and two-point LMTD overestimates its area
by 129--223\,\% at $\varepsilon = 0$. The thesis also shows that the heater
inlet temperature, raised by recuperation, is the floor on heat recovery.
\item \textbf{A fixed-source figure of merit.} Net work per kilogram of
waste gas (Eq.~\ref{eq:wgas}), combined with cost at the delivered output,
is shown to rank cycles differently from thermal efficiency, and to be the
quantity on which the only CO$_2$ economic win (tCO$_2$ at
200\,\textdegree C) rests.
\item \textbf{A selection framework for Australian industry} validated
against 10,000-sample Monte Carlo uncertainty and transparent about its
revisions from the interim stage.
\end{enumerate}

\section{Recommendations for Future Work}
\label{sec:future_work}

\begin{enumerate}[leftmargin=1.8em,itemsep=4pt]
\item \textbf{Cascade and split-flow sCO$_2$ layouts} for deep recovery,
evaluated on the same fixed-source basis, since the simple recuperated
layout's shallow recovery is its principal economic handicap.
\item \textbf{Joint optimisation of CO$_2$ pressures and recuperator
effectiveness for $w_{gas}$}, as was done for the steam pressure here,
with the chiller or cooling-tower system costed explicitly as a function of
site wet-bulb temperature.
\item \textbf{Vendor-based costs for sub-5\,MW CO$_2$ turbomachinery and
PCHEs}, including pressure drop, to replace the extrapolated Weiland
correlations; and a condensing-turbine cost for steam to replace the
back-pressure benchmark.
\item \textbf{Dynamic simulation} with turbomachinery maps and heater
re-integration at part load, to replace the single-coefficient Stodola model.
\item \textbf{Identification of Australian sites with natural cold sinks}
(cold groundwater, deep-water intakes, LNG regasification cold) where the
tCO$_2$ cycle's recovery advantage could be realised.
\item \textbf{Lifecycle greenhouse-gas and carbon-price sensitivity}, which
would favour the cycle that recovers the most energy rather than the one
with the lowest capital cost.
\end{enumerate}

\section{Concluding Statement}
\label{sec:concluding}

Industrial waste heat recovery in Australia is a sound investment above
300\,\textdegree C: at 2026 costs and prices a steam Rankine plant at a
cement preheater or an electric arc furnace returns its capital in three to
four years and is profitable in essentially every scenario examined. The
CO$_2$ cycles do not displace it. Their compactness is real and their
part-load behaviour is better, but at Australian sink temperatures the
supercritical Brayton cycle loses the near-critical compression on which its
efficiency depends, and the transcritical Rankine cycle --- which recovers
more energy than steam from any stream below 400\,\textdegree C --- needs a
condenser colder than the country can supply without spending more on
refrigeration than the cycle earns. The thermodynamics that make CO$_2$
attractive in the literature are not wrong; they belong to cold climates,
hot sources and optimised layouts. For an Australian operator in 2026 the
recommendation is steam at 350\,\textdegree C and above, a careful site-
specific study at 200\,\textdegree C, and transcritical CO$_2$ wherever a
natural cold sink makes its deep recovery usable.
